\exercisesection{\S~5}

\begin{enumerate}
\def\labelenumi{\arabic{enumi})}
\item
  证明第一章 §5 习题 6 中所考虑的可解李代数不同构于任何可分解李代数。
\item
  设 u（相应地 v）是 V 的非零半单（相应地幂零）自同态。则映射 \(\lambda u \mapsto \lambda v (\lambda \in k)\) 是从 g = ku 到 \(g' = kv\) 的同构，但它不把半单元素映为半单元素。
\item
  设 \(u\) 是 \(V\) 的既非半单亦非幂零的自同态。则 \(\mathfrak{g} = ku\) 不可分解，但 \(\operatorname{ad}_{\mathfrak{g}}\mathfrak{g}\) 可分解。
\end{enumerate}

¶ 4) 设 g 是 gl(V) 的可分解李子代数。设 q 是 g 的李子代数，其恒同表示半单。存在 \(a \in \operatorname{Aut}_{e}(\mathfrak{g})\) 使 \(a(\mathfrak{q})\) 含于命题 7 的子代数 m。（模仿命题 7 (ii) 的证明。）

¶ 5) 设 g 是 gl(V) 的李子代数。若对一切 \(x \in \mathfrak{g}\) ，\(x\) 的复本（第一章 §5 习题 14）都属于 g，则称 g 为代数的。这样的代数是可分解的。

\begin{enumerate}
\def\labelenumi{\alph{enumi})}
\tightlist
\item
  用 \(a(\mathfrak{g})\) 表示 \(\mathfrak{gl}(\mathrm{V})\) 中包含 \(\mathfrak{g}\) 的最小代数子代数。则
\end{enumerate}

\[
a (\mathfrak {g}) \supset e (\mathfrak {g}) \supset \mathfrak {g}.
\]

给出一个 \(a(\mathfrak{g})\) 与 \(e(\mathfrak{g})\) 相异的例子（取 V 为 2 维，\(\mathfrak{g}\) 为 1 维）。

\begin{enumerate}
\def\labelenumi{\alph{enumi})}
\setcounter{enumi}{1}
\item
  证明：若 \(\mathfrak{n}\) 是 \(\mathfrak{g}\) 的理想，则 \(\mathfrak{n}\) 与 \(a(\mathfrak{n})\) 都是 \(a(\mathfrak{g})\) 的理想，且 \([a(\mathfrak{g}), a(\mathfrak{n})] = [\mathfrak{g}, \mathfrak{n}]\) （模仿命题 4 的证明）。推出 \(\mathcal{D}^i a(\mathfrak{g}) = \mathcal{D}^i \mathfrak{g}\) 对 \(i \geq 1\) 成立，且 \(\mathcal{C}^i a(\mathfrak{g}) = \mathcal{C}^i \mathfrak{g}\) 对 \(i \geq 2\) 成立。
\item
  证明由幂零元素组成的李代数都是代数的。\(^{8}\)
\end{enumerate}

¶ 6) 设 g 是 \(\mathfrak{gl}(V)\) 的半单子代数，\(\mathbf{T}(V) = \bigoplus_{n=0}^{\infty} \mathbf{T}^{n}(V)\) 是 V 的张量代数，\(\mathbf{T}(V)^{\mathfrak{g}}\) 是 \(\mathbf{T}(V)\) 中在 g 下不变的元素之集（参见第三章附录），\(\overline{g}\) 表示这样的 \(u \in \mathfrak{gl}(V)\) 组成的集合：u.x = 0 对一切 \(x \in \mathbf{T}(V)^{\mathfrak{g}}\) 成立。我们将证明 \(\overline{g} = g\) 。

\begin{enumerate}
\def\labelenumi{\alph{enumi})}
\item
  证明 g 在 V 的对偶 V* 上的表示同构于它在 \(\bigwedge^{p-1}V\) 上的表示，其中 \(p = \dim V\) （利用 g 含于 \(\mathfrak{sl}(V)\) 这一事实）。推出 \(\mathbf{T}_{n,m} = \mathbf{T}^{n}(\mathrm{V}) \otimes \mathbf{T}^{m}(\mathrm{V}^{*})\) 中在 g 下不变的每个元素都在 \(\overline{g}\) 下不变，且 \(\overline{g}\) 是代数的（习题 5）。
\item
  设 W 是某个 \(T_{n,m}\) 的向量子空间。假设 W 在 g 下稳定。证明 W 在 \(\overline{g}\) 下稳定（若 \(e_{1},\ldots,e_{r}\) 是 W 的基，注意 \(e_{1}\wedge\cdots\wedge e_{r}\) 在 g 下不变，从而在 \(\overline{g}\) 下不变）。
\end{enumerate}

推出 g 是 \(\overline{g}\) 的理想，且 \(\overline{g}/g\) 交换（参见命题 4 的证明）。我们有 \(\overline{g} = g \times c\) ，其中 c 是 \(\overline{g}\) 的中心。

\begin{enumerate}
\def\labelenumi{\alph{enumi})}
\setcounter{enumi}{2}
\item
  设 R 是 \(\mathfrak{gl}(V)\) 中由 1 与 g 生成的结合子代数。证明 c 含于 R 的中心（注意 \(\overline{g}\) 含于 R 的双交换子，而后者等于 R）。推出 c 的元素都是半单的。
\item
  设 \(x \in c\) 。证明 x 的复本属于 c（第一章 §5 习题 14）。证明 \(\operatorname{Tr}(sx) = 0\) 对一切 \(s \in g\) 成立；推出 \(\operatorname{Tr}(sx) = 0\) 对一切 \(s \in \overline{g}\) 成立，从而 x 幂零（同前引）。
\item
  证明 \(\mathfrak{c} = 0\) 且 \(\overline{\mathfrak{g}} = \mathfrak{g}\) ，把 \(c\) 与 \(d\) ) 结合起来。
\end{enumerate}

\begin{enumerate}
\def\labelenumi{\arabic{enumi})}
\setcounter{enumi}{6}
\tightlist
\item
  设 \(\mathfrak{g}\) 是 \(\mathfrak{gl}(\mathrm{V})\) 的李子代数。设 \(m, n\) 是两个整数 \(\geq 0\) ，W 与 \(\mathrm{W}'\) 是 \(\mathbf{T}^m(\mathrm{V}) \otimes \mathbf{T}^n(\mathrm{V}^*)\) 的两个向量子空间，其中 \(\mathrm{V}^*\) 是 V 的对偶。假设 \(\mathrm{W}' \subset \mathrm{W}\) ，且 W 与 \(\mathrm{W}'\) 在 \(\mathfrak{g}\) 于 \(\mathbf{T}^m(\mathrm{V}) \otimes \mathbf{T}^n(\mathrm{V}^*)\) 上的自然表示下稳定。证明 W 与 \(\mathrm{W}'\) 此时在 \(e(\mathfrak{g})\) 下稳定。若用 \(\pi\) 表示 \(e(\mathfrak{g})\) 由此在 \(\mathrm{W/W}'\) 上给出的表示，证明 \(\pi e(\mathfrak{g})\) 是 \(\pi(\mathfrak{g})\) 的可分解包络（利用定理 1）。
\end{enumerate}

推出 \(\operatorname{ad} e(\mathfrak{g})\) 是 \(\operatorname{ad} \mathfrak{g}\) 在 \(\mathfrak{gl}(\mathfrak{g})\) 中的可分解包络。

\begin{enumerate}
\def\labelenumi{\arabic{enumi})}
\setcounter{enumi}{7}
\tightlist
\item
  设 g 是 \(\mathfrak{gl}(V)\) 的李子代数，h 是 g 的嘉当子代数。
\end{enumerate}

\begin{enumerate}
\def\labelenumi{\alph{enumi})}
\tightlist
\item
  证明 \(e(\mathfrak{g}) = e(\mathfrak{h}) + \mathcal{D}\mathfrak{g} = e(\mathfrak{h}) + \mathfrak{g}\) 。
\end{enumerate}

（注意 \(e(\mathfrak{h}) + \mathcal{D}\mathfrak{g}\) 可分解（定理 1 的推论 1），包含 \(\mathfrak{g} = \mathfrak{h} + \mathcal{D}\mathfrak{g}\) ，且含于 \(e(\mathfrak{g})\) ；故它就是 \(e(\mathfrak{g})\) 。）

\begin{enumerate}
\def\labelenumi{\alph{enumi})}
\setcounter{enumi}{1}
\item
  我们有 \(e(\mathfrak{h}) \cap \mathfrak{g} = \mathfrak{h}\) （注意 \(e(\mathfrak{h}) \cap \mathfrak{g}\) 幂零）。
\item
  设 \(x\) 是 \(e(\mathfrak{h})\) 在 \(e(\mathfrak{g})\) 中的正规化子的元素。证明 \(x \in e(\mathfrak{h})\) 。（写 \(x = y + z\) ，其中 \(y \in e(\mathfrak{h}), z \in \mathfrak{g}\) ，参见 a)；注意 \([z, \mathfrak{h}] \subset e(\mathfrak{h}) \cap \mathfrak{g} = \mathfrak{h}\) ，故 \(z \in \mathfrak{h}\) 。）
\end{enumerate}

\(d\) ) 证明 \(e(\mathfrak{h})\) 是 \(e(\mathfrak{g})\) 的嘉当子代数。

\begin{enumerate}
\def\labelenumi{\arabic{enumi})}
\setcounter{enumi}{8}
\tightlist
\item
  设 \(\mathfrak{g}\) 是 \(\mathfrak{gl}(\mathrm{V})\) 的李子代数。证明 §3 习题 16 的条件 (i)、(ii)、(iii) 等价于：
\end{enumerate}

\begin{enumerate}
\def\labelenumi{(\alph{enumi})}
\setcounter{enumi}{21}
\tightlist
\item
  \(\mathfrak{g}\) 可分解，且与 \(\mathfrak{g} / \mathfrak{n}_{\mathrm{V}}(\mathfrak{g})\) 有相同的秩。
\end{enumerate}

（若 \(\mathfrak{h}\) 是 \(\mathfrak{g}\) 的嘉当子代数，则条件“\(\mathfrak{g}\) 与 \(\mathfrak{g}/\mathfrak{n}_{\mathrm{V}}(\mathfrak{g})\) 有相同的秩”等价于说 \(\mathfrak{h} \cap \mathfrak{n}_{\mathrm{V}}(\mathfrak{g}) = 0\) ，即 \(\mathfrak{h}\) 没有 \(\neq 0\) 的幂零元素（参见 §2 习题 14）。推出 (i) 与 (v) 等价。）

\begin{enumerate}
\def\labelenumi{\arabic{enumi})}
\setcounter{enumi}{9}
\tightlist
\item
  设 \(k'\) 是 \(k\) 的一个扩张，\(\mathfrak{g}'\) 是下式的 \(k'\) -李子代数
\end{enumerate}

\[
\mathfrak {g l} (\mathrm{V} \otimes_ {k} k ^ {\prime}) = \mathfrak {g l} (\mathrm{V}) \otimes_ {k} k ^ {\prime}.
\]

\begin{enumerate}
\def\labelenumi{\alph{enumi})}
\item
  证明存在 \(\mathfrak{gl}(V)\) 的最小李子代数 g，使得 \(g\otimes_{k}k'\) 包含 \(g'\) 。
\item
  假设 \(k'\) 代数闭，并用 G 表示 \(k'\) 的 k-自同构群；此群以自然方式作用于 \(V \otimes_{k} k'\) 。证明 \(g \otimes_{k} k'\) 是由 \(g'\) 在 G 下的共轭元生成的李子代数（利用 k 是 \(k'\) 中 G-不变元所成的域这一事实）。
\item
  证明 \(\mathfrak{g}\) 可分解，若 \(\mathfrak{g}'\) 可分解。（化归到 \(k'\) 代数闭的情形，并利用 \(b\) ) 以及定理 1 的推论 1 与 3。）
\end{enumerate}

¶ 11) 作为例外，本习题中我们假设 k 是特征为 p \textgreater{} 0 的完满域。

设 g 是李 p-代数（第一章 §1 习题 20）。若 \(x \in g\) ，用 \(\langle x \rangle\) 表示包含 x 的最小李 p-子代数。它是交换的，且作为 k-向量空间由诸 \(x^{p^{i}}\) 生成，其中 \(i = 0, 1, \ldots\) 。称 x 为幂零的（相应地半单的），若 \(\langle x \rangle\) 的 p-映射幂零（相应地双射）。

\begin{enumerate}
\def\labelenumi{\alph{enumi})}
\item
  证明 x 可以唯一地分解为形式 \(x = s + n\) ，其中 \(s, n \in \langle x \rangle\) ，s 半单，n 幂零（应用第一章 §1 习题 23）。若 f 是从 g 到 \(\mathfrak{gl}(\mathrm{V})\) 的 p-同态，则 \(f(s)\) 与 \(f(n)\) 是自同态 \(f(x)\) 的半单分量与幂零分量；这特别适用于 f = ad。
\item
  g 的子代数称为可分解的，若它包含其元素的半单分量与幂零分量。证明：若 b 与 c 是 g 的向量子空间且 \(b \subset c\) ，则 \(x \in g\) 中使得 \([x, c] \subset b\) 的那些 x 组成的集合是可分解的（与命题 3 的证明相同）；特别地，g 的每个嘉当子代数都可分解。
\item
  设 t 是 g 的由半单元素组成的交换子代数，且对此性质极大。设 h 是 t 在 g 中的交换子。设 \(x \in h\) ，\(x = s + n\) 是其典则分解；由于 h 可分解（参见 b)），\(s, n \in h\) 。证明由 t 与 s 生成的子代数交换且由半单元素组成，故与 t 一致。推出 \(ad_{h}x = ad_{h}n\) 幂零，从而 h 幂零。由于 \(\mathfrak{h} = \mathfrak{g}^{0}(\mathfrak{h})\) ，h 是 g 的嘉当子代数（§2 命题 4）。
\end{enumerate}

特别地，有限域上的每个李 p-代数都有嘉当子代数。\(^{9}\)

